\section{Natural-language deployment: who holds the signing key}
\label{sec:nl-deployment}

A user can already describe what they want running in a sentence and a short paragraph
of configuration, and something on \Flux will turn that description into a running
application. What differs sharply, and matters to anyone deciding how much to trust the
result, is which of two \emph{different systems} turns the description into a signed,
on-chain transaction. One is the protocol's own path: the string to be signed goes to a
device the user holds, the user's key signs it there, and a fully compromised assistant
is limited to proposing something bad. The other is a third-party operator's service:
the operator holds the key, the operator signs, and what stands between an autonomous
caller and a live deployment is a bearer token that authenticates the caller to the party
already holding the key (\S\ref{sec:paying:facilitator}). These are not two modes of one
product, and this section does not present them as a choice the protocol offers — the
protocol offers exactly one of them and is indifferent to the existence of the other
(Property~\ref{prop:chainindifferent}). The section states each precisely, gives each its
own security property rather than one property for ``the system,'' and argues that for
the protocol's path the open problem is not custody at all — it is whether the artifact
placed in front of the signer can actually be understood before it is signed.

\subsection{What ships today}

Three pieces of tooling turn a natural-language or near-natural-language description into
a deployment. All three are \shipped.

\textbf{Orbit.} A wizard (React front end plus an Express server, publicly reachable at
\texttt{orbit.\allowbreak{}runonflux.\allowbreak{}com} per its own configuration) walks a user through a Git
repository URL, a resource plan and a small number of deployment options, then emits a
\specv{8} application specification whose single compose component runs the shared image
\texttt{runonflux/\allowbreak{}orbit:latest} \shipped
\srccode{deploy-with-git-frontend/\allowbreak{}config/\allowbreak{}defaults.js}{3}
\srccode{deploy-with-git-frontend/\allowbreak{}src/\allowbreak{}services/\allowbreak{}deployService.js}{453--490}.
That specification is not a pointer to a build the Flux team performs on the user's
behalf: once FluxOS places it on a node, the container itself clones the target
repository, auto-detects its language and framework by an ordered set of file-presence
checks, installs dependencies, builds and runs it, and continues to manage releases,
rollback and webhook-triggered redeploys from inside the running instance \shipped
\srccode{orbit/\allowbreak{}README.md}{5,20} \srccode{orbit/\allowbreak{}entrypoint.sh}{6--9}
\srccode{orbit/\allowbreak{}lib/\allowbreak{}build.sh}{304--431}. Orbit is not a separate build farm; the
container \emph{is} the workload, running on whichever node FluxOS happens to place that
application instance on \shipped \srccode{orbit/\allowbreak{}entrypoint.sh}{1--9}
\srccode{orbit/\allowbreak{}README.md}{39--57}.

\textbf{Deploy-from-Git.} The end-to-end path from a repository to a running application
is: a user or agent supplies a Git URL and configuration; the frontend server analyzes
the repository (provider, public/private access, branches, a detected port, monorepo
layout, an imported \texttt{flux.json}\slash\texttt{flux.yaml}) \shipped
\srccode{deploy-with-git-frontend/\allowbreak{}server/\allowbreak{}orbit/\allowbreak{}repositoryService.js}{16--53}; a \specv{8}
specification is built around the Orbit image \shipped
\srccode{deploy-with-git-frontend/\allowbreak{}src/\allowbreak{}services/\allowbreak{}deployService.js}{453--489}; the
specification is sent to Flux's own \texttt{POST /\allowbreak{}apps/\allowbreak{}verifyap\allowbreak{}pregistr\allowbreak{}ationspe\allowbreak{}cificati\allowbreak{}ons}
for normalization \shipped
\srccode{deploy-with-git-frontend/\allowbreak{}src/\allowbreak{}services/\allowbreak{}deployService.js}{521--539}; it is signed
(mechanism below); and the signed transaction is submitted directly to
\texttt{api.\allowbreak{}runonflux.\allowbreak{}io} via \texttt{POST /\allowbreak{}apps/\allowbreak{}appregister} \shipped
\srccode{deploy-with-git-frontend/\allowbreak{}src/\allowbreak{}services/\allowbreak{}deployService.js}{547--565}. Neither the
frontend server nor Orbit itself ever clones or builds the user's code before that point
— that work happens later, inside the Orbit container, once placement has occurred.

\textbf{\texttt{DeploymentYAML}.} A gRPC method,
\texttt{DeploymentYAML(YAMLPrompt) $\to$ YAMLResponse}, generates Kubernetes YAML from a
natural-language prompt \shipped
\srccode{pouwbackend/\allowbreak{}grpc\_server/\allowbreak{}market\_place/\allowbreak{}deployment.go}{1784--1812}. The
generator first vector-searches existing marketplace application YAMLs for
retrieval-augmented context \shipped
\srccode{pouwbackend/\allowbreak{}ai/\allowbreak{}fluxai.go}{188--191,241--244}, then issues a chat-completion
request against a hard-coded system preamble instructing the model to emit only raw YAML,
with the model pinned to ``Llama 3.1 8B'' and dispatched either to FluxAI's own
\texttt{/\allowbreak{}v1/\allowbreak{}chat/\allowbreak{}completions} endpoint or to Google AI \shipped
\srccode{pouwbackend/\allowbreak{}ai/\allowbreak{}yaml.go}{30--41} \srccode{pouwbackend/\allowbreak{}ai/\allowbreak{}fluxai.go}{101,116,212,219--225,266,278--284}.
Natural-language deployment, in other words, is not hypothetical; a form of it ships
today as a retrieval-augmented, LLM-backed generator. What the extraction did
\emph{not} find is a code path chaining that generated YAML directly into
\texttt{/\allowbreak{}deployment/\allowbreak{}start} — whether the output is ever auto-submitted or always
requires a human or agent to copy it forward is unresolved from the evidence available,
and the claim made here is limited to generation, not to confirmed unattended
submission \srcinferred.

\subsection{Two signing paths, belonging to two systems}
\label{sec:nl:paths}

The three pieces of tooling above all converge on one signing step, and that step does
not have one answer — it depends on how the caller authenticated, and therefore on which
system produces the signature.

\begin{definition}[Local signing and operator signing]
\label{def:nl:paths}
A \emph{local-signing path} is one on which the string to be signed is built by a
frontend and relayed to a device that holds the private key, with the signature
computed and returned by that device. An \emph{operator-signing path} is one on which a
remote service that itself holds the private key computes the signature and returns it,
authorized by a bearer credential that authenticates the caller to that service. The word
\emph{delegated} does not fit the second and this paper does not use it as a description
of authority there: delegation presupposes a principal holding the authority being
delegated, and on this path the user never holds the key
(\S\ref{sec:paying:facilitator}, Property~\ref{prop:custody},
Corollary~\ref{cor:nodelegation}).
\end{definition}

The first is a path of the protocol: \Flux verifies a signature against an address and
takes no view of where the corresponding key sits, so a key on a user's own device is
exactly what the network expects and requires no service anywhere. The second is a path
of a company's product; the network cannot see it and does not depend on it
(Property~\ref{prop:chainindifferent}). The remainder of this subsection states each
separately, and it does not average them.

On the local-signing path, a user connected through the \SSP wallet extension or through
ZelCore has the frontend build the string-to-sign and relay it to \texttt{window.ssp.request('sign', \ldots)}
or to a ZelCore deep link plus a WebSocket channel; the private key never leaves the
user's own device, and the frontend only ever sees the resulting signature \shipped
\srccode{deploy-with-git-frontend/\allowbreak{}src/\allowbreak{}services/\allowbreak{}deployService.js}{697--706,717--772}.
Here the trust boundary really is the signing step: a fully compromised assistant can
construct and propose an arbitrary specification, but it cannot produce a valid signature
over it, and therefore cannot deploy or spend. That is the property this section's opening
claims for the protocol's path, and here it holds.

On the operator-signing path — a \FluxCore SSO login authenticated by Firebase — the
frontend server does not hold or see a Flux private key at any point. It POSTs the message
and the user's Firebase ID token to a remote service,
\texttt{https:/\allowbreak{}/\allowbreak{}service.\allowbreak{}fluxcore.\allowbreak{}ai/\allowbreak{}api/\allowbreak{}signMessage},
and receives a signature back; the service signs \emph{on the user's behalf} \shipped
\srccode{deploy-with-git-frontend/\allowbreak{}src/\allowbreak{}services/\allowbreak{}deployService.js}{677--691}
\srccode{deploy-with-git-frontend/\allowbreak{}server/\allowbreak{}flux/\allowbreak{}fluxSession.js}{60--73}. That remote service
is the third-party payment facilitator of \S\ref{sec:paying:facilitator}: the same
architecture, in which the operator generates and holds both of a user's private keys and
the user never sees either \srcintent{design intake}
\srccode{console-api/\allowbreak{}src/\allowbreak{}db/\allowbreak{}models/\allowbreak{}FluxIdentity.ts}{13--28}; the identification of the two
services follows from the observed protocol rather than from a read of this signer, which
is not in the corpus \srcinferred. This is the path
every server-side call uses, including registration
(\texttt{DeploymentService.\allowbreak{}register()} calling \texttt{sessions.\allowbreak{}signMessage(session, message)},
authenticated only by the caller's bearer token) and updates
(\texttt{UpdateService.\allowbreak{}submit()}) \shipped
\srccode{deploy-with-git-frontend/\allowbreak{}server/\allowbreak{}orbit/\allowbreak{}deploymentService.js}{158--173}
\srccode{deploy-with-git-frontend/\allowbreak{}server/\allowbreak{}orbit/\allowbreak{}updateService.js}{68--79}. It is also the
path the shipped Model Context Protocol (MCP) tools \texttt{deploy\_app},
\texttt{update\_app} and \texttt{renew\_app} use \shipped
\srccode{deploy-with-git-frontend/\allowbreak{}server/\allowbreak{}mcp/\allowbreak{}createServer.js}{13--53,66--266}
\srccode{deploy-with-git-frontend/\allowbreak{}server/\allowbreak{}mcp/\allowbreak{}router.js}{13--44}. The \texttt{deploy\_app}
tool's own description states the mechanism plainly: ``Revalidate, Firebase-sign,
register, and test-install an Orbit app'' \shipped
\srccode{deploy-with-git-frontend/\allowbreak{}server/\allowbreak{}mcp/\allowbreak{}createServer.js}{174--177}. These tools, along
with \texttt{control\_instance} and \texttt{trigger\_build}, are marked
\texttt{destructive\allowbreak{}Hint:\allowbreak{}true} and execute against live Flux state given only a valid
bearer token \shipped \srccode{deploy-with-git-frontend/\allowbreak{}server/\allowbreak{}mcp/\allowbreak{}createServer.js}{170--266}.
Given such a token, an agent can construct, sign and submit a deployment end to end with
no human in the loop and no separate signing step \shipped
\srccode{deploy-with-git-frontend/\allowbreak{}server/\allowbreak{}mcp/\allowbreak{}createServer.js}{66--266}
\srccode{deploy-with-git-frontend/\allowbreak{}server/\allowbreak{}orbit/\allowbreak{}services.js}{91--104}.

\begin{table}[H]
\centering
\begin{tabular}{p{2.7cm}p{4.6cm}p{5.2cm}}
\hline
\textbf{Path, and whose it is} & \textbf{Who produces the signature} & \textbf{Bound on a fully compromised signing credential} \\
\hline
Local signing (\SSP, ZelCore)\newline\emph{the protocol's path} &
The user's own device; the frontend only builds the string-to-sign and relays the
resulting signature. \shipped\newline
\srccode{deploy-with-git-frontend/\allowbreak{}src/\allowbreak{}services/\allowbreak{}deployService.js}{697--706,717--772} &
A compromised assistant can propose a bad specification; it cannot deploy or spend,
because it never holds a key capable of producing a valid signature over it. \shipped\newline
\srccode{deploy-with-git-frontend/\allowbreak{}src/\allowbreak{}services/\allowbreak{}deployService.js}{697--706,717--772} \\
\hline
Operator signing (\FluxCore SSO)\newline\emph{a third party's service, built on \Flux} &
A remote service, \texttt{service.\allowbreak{}fluxcore.\allowbreak{}ai/\allowbreak{}api/\allowbreak{}signMessage}, which holds the key and
is authorized only by a Firebase bearer token. \shipped\newline
\srccode{deploy-with-git-frontend/\allowbreak{}src/\allowbreak{}services/\allowbreak{}deployService.js}{677--691}
\srccode{deploy-with-git-frontend/\allowbreak{}server/\allowbreak{}flux/\allowbreak{}fluxSession.js}{60--73} &
Nothing: no per-transaction cap, no rate limit, no budget, no scope, and no revocation
beyond invalidating the token itself. \srcinferred\newline
\srcintent{mechanism specification} \\
\hline
\end{tabular}
\caption{The bound on a fully compromised signing credential, by path. The two rows are
properties of two different systems and do not compose into a property of ``the system''.}
\end{table}

Non-custody is a property of the protocol's path. It is not a property of the operator's
service, and the operator's service is not part of \Flux. Stating the claim at the level
of ``the assistant configures, the user signs, therefore a compromised assistant cannot
spend'' would be false as a description of the operator path, because there an assistant
holding a Firebase bearer token \emph{is} the signer, indirectly, through a service that
will sign whatever it is asked to sign.

Two consequences, and the second is the one readers get wrong. First, the bound on that
token today is nothing at all, and nothing in the protocol could have made it otherwise:
the token is not a delegation and confers no authority — it authenticates a caller to the
party that already holds the key (\S\ref{sec:paying:facilitator},
Property~\ref{prop:custody}). Second, and therefore, \textbf{\S17's construction does not
apply to this path.} A delegation certificate bounds an authority a principal holds and
grants; here the principal holds nothing, so there is no grant to bound and the
certificate has nowhere to attach (Corollary~\ref{cor:nodelegation}). What would bound the
operator path is operator-side policy — rate limits, per-token caps, an approval step —
which is a product decision belonging to the operator, and this paper specifies none of it
and does not track it. \S17 is written for the other case: a principal \(\principal\) who
holds a key on the protocol's path and grants an agent \(\agent\) a scope \(\scopeset\)
under a budget \(\budget\) — where \S17 also states why enforcing \(\budget\) on the
\Flux L1 as deployed is not currently possible without a new trusted party or a consensus
change.

The operator may choose to hand users the keys it holds for them — \(K_u\) in the notation
of \S\ref{sec:paying:facilitator} — at which point its path becomes a local-signing one;
it states that it may
\srcintent{design intake}. That is the operator's decision and not a
protocol change — no consensus rule, no \FluxOS release and no \appspec field moves either
way — and this paper records it without tracking it
(\S\ref{sec:paying:facilitator}, Remark~\ref{rem:operatoroption}).

\subsection{Reviewability: the question that remains}

On the protocol's path the security boundary is exactly the signing step, and
that boundary is only as strong as the signer's ability to understand what is being
signed. Everything in this subsection is about that path; on the operator's service the
question does not arise, because the party that reviews and the party that signs are the
same party and it is not the user. For that review to be a meaningful check rather than a formality, the interface
would need to surface, at minimum: the total cost and duration of the deployment; every
port the specification exposes; the provenance of every image \emph{by content digest,
not by mutable tag}; the resource commitments (cpu, memory, storage) the specification
requests; the data-retention consequences of eviction or non-renewal (\S16); and any
node-targeting or geolocation constraint the specification carries. \planned. These are
not arbitrary — they are structurally the same fields \S17's delegation scope
\(\scopeset = (\mathcal{N}, \resvec^{\max}, \ninst^{\max}, \mathcal{I}, \dur^{\max}, \mathcal{G})\)
already names, because a human reviewing a specification before signing it and a
delegation certificate bounding an agent's authority to construct one are answering the
same question. \planned \srcintent{mechanism specification}.

The concrete failure this addresses is not hypothetical. The deployed tooling emits
\texttt{runonflux/\allowbreak{}orbit:latest} — a mutable tag, not a content digest — with nothing in
the specification-construction code pinning a \texttt{sha256:} digest \shipped
\srccode{deploy-with-git-frontend/\allowbreak{}src/\allowbreak{}services/\allowbreak{}deployService.js}{453--490}. A user
signing that specification is not reviewing the code that will run; they are reviewing
a name that resolves to whatever the tag currently points to at the moment a Flux node
pulls it. The review surface itself compounds this: the wizard's final confirmation step
renders a reconstructed summary object — name, ports, domain, polling interval, runtime,
geolocation, environment-variable \emph{keys} only — and not the literal \specv{8}
payload that gets signed, which omits the image reference, owner, contacts, expiry and
the full environment-variable array present in the actual signed string
\shipped \srccode{deploy-with-git-frontend/\allowbreak{}src/\allowbreak{}components/\allowbreak{}wizard/\allowbreak{}Step4Review.jsx}{280--338}.
The literal artifact signed is only ever visible as the \texttt{dataToSign} string passed
to the signer, and nothing in the extraction found it rendered to the user anywhere
\shipped \srccode{deploy-with-git-frontend/\allowbreak{}src/\allowbreak{}services/\allowbreak{}deployService.js}{667--670}. An
unreviewable specification presented for signature reduces the security boundary to
theatre: the signing step still cryptographically binds the signer to \emph{something},
but not demonstrably to the thing the signer believes they approved.

It is worth noting, briefly, that the draft \specv{9} application specification already
carries several of the fields a meaningful review would need — a required top-level
\texttt{duration} field, per-component \texttt{cpu}\slash\texttt{memory}\slash\texttt{rootFsGb}, and
an explicit \texttt{ports} map — none of which existed as required fields in \specv{8}
\inprogress \srccode{flux-spec/\allowbreak{}schema/\allowbreak{}v9.json}{62--66,103--117,129--150}. What that draft
does not yet specify, as of this writing, is any signature field at all: its
canonicalization and fingerprint recipe is described only as ``recommended,'' not
enforced, and no signature or signer-identity field appears anywhere in the schema
\inprogress \srccode{flux-spec/\allowbreak{}validation/\allowbreak{}canonicalization.md}{53--63}. The reviewability
question this section raises is therefore not one \specv{9} answers by construction; it
is a gap the next specification version inherits along with the fields it would need to
close it.

\begin{property}
The signing step is a meaningful trust boundary if and only if the signer can, from the
artifact presented for signature alone, determine (i) the set of programs that may
execute as a result of that signature, and (ii) the maximum value that may be spent as a
result of it.
\end{property}

The deployed interface fails condition (i) outright: the image reference a signed
specification authorizes is a mutable tag rather than a fixed digest, so two signatures
over ``the same'' specification, at different times, can authorize different code
\shipped \srccode{deploy-with-git-frontend/\allowbreak{}src/\allowbreak{}services/\allowbreak{}deployService.js}{453--490}, and
the object rendered for review is a reconstruction rather than the literal payload being
signed \shipped \srccode{deploy-with-git-frontend/\allowbreak{}src/\allowbreak{}components/\allowbreak{}wizard/\allowbreak{}Step4Review.jsx}{280--338}.
Condition (ii) fares somewhat better in structure — plan, resource ceiling and billing
period are concrete fields the tooling computes before submission, including an explicit
blocks-from-months expiry calculation \shipped
\srccode{deploy-with-git-frontend/\allowbreak{}src/\allowbreak{}services/\allowbreak{}deployService.js}{102,108--110} — but
because the rendered review is, again, a reconstruction, a signer confirms the wizard's
arithmetic rather than reading the number that will actually be signed. Digest pinning
is the cheapest available remedy for condition (i), because it requires no protocol
change, no consensus change and no new specification field: it only requires
substituting a resolved digest for a mutable tag at the point where the tooling already
constructs the specification, which is a change confined entirely to the tooling that
already runs today.

\subsection{Self-hosting}

The tooling that performs natural-language deployment does not sit outside the network
it targets. Orbit is deployed as an ordinary \Flux application and performs its own
clone/build/run cycle inside that application's own container, on whichever node
FluxOS happens to place it — it is explicitly not a separate, team-operated build farm
\shipped \srccode{orbit/\allowbreak{}README.md}{5,20} \srccode{orbit/\allowbreak{}entrypoint.sh}{1--9}. The
language model behind \texttt{DeploymentYAML} is served from FluxAI's own inference
capacity, and that capacity is itself deployed as workloads onto \FluxEdge nodes through
the same GitOps pipeline used for any other model \shipped
\srccode{fluxai-fluxedge/\allowbreak{}custom-workflow-fluxone.yaml}{1--52}. An assistant that deploys
applications on \Flux is, in this narrow sense, expected to run as an ordinary tenant of
the network rather than as privileged infrastructure. That is a pleasing property, worth
stating once — but it is not a security property, and it should not be dressed up as
one. Nothing about running as an ordinary tenant bounds what a compromised
operator-signing credential can do; the bound in \S\ref{sec:nl:paths} is unaffected by
where the assistant itself happens to run.

\subsection{Toward bounded automation}

A principal \emph{who holds a key} and explicitly delegates under \S17's certificate
construction would obtain automation whose worst case is bounded rather than open-ended:
the certificate fixes a scope \(\scopeset\) — permitted application names, a resource
ceiling, a digest-pinned image set, a maximum duration and permitted geolocation
constraints — and, where enforceable, a spending budget \(\budget\). \planned
\srcintent{mechanism specification}. The qualifier is load-bearing and is
the reason this paragraph sits under the protocol's path: the construction presupposes a
principal-held key, so it is available to a user on the local-signing path and to
autonomous software holding its own key, and it is not available on the operator-signing
path at all (Corollary~\ref{cor:nodelegation}). This section does not restate that
construction; it only notes the relationship — that construction is what would let a
natural-language deployment flow run end to end without the two open failures described
above, an unpinned image and an unbounded bearer token, both being in force at once.

\subsection{What would make the boundary real}

Three changes, taken together, would make the trust boundary on the protocol's path
correspond to what it is presented as being — and, incidentally, would give an operator
that chose to bound its own service something worth enforcing: digest pinning by default,
so that \(\mathcal{I}\)
names fixed content rather than a mutable tag; a specification diff presented before
signature, so the artifact a signer reviews is the literal payload being signed rather
than a reconstructed summary of it; and a cost estimate the signer can check against the
specification's own fields rather than trust from a wizard screen. \planned \srcinferred.
None of the three requires a consensus change or a new trusted party; all three are
changes to tooling that already runs today.

%% Drain this section's float queue before the next begins.
\FloatBarrier
